% Pursuing Stacks (cotes 134-2 to 134-8): archive cote 134-6, site page 75 (PDF page 76), typescript page 393.
% typescript. Draft compared: ps5.tex:5233-5298.
% AI-transcribed by Claude Opus 5.5 (claude-opus-5-5, adaptive thinking, effort high), Message Batches API, batch ps3 (msgbatch_01XdFmuVVBAcgc7EFxdtCqsj), 2026-11-12; not yet checked by eye.
\scanpage{134-6/75}
\mspage{393}
an arbitrary \hand{$\mathcal{A}$} satisfying our assumptions, even when \hand{$\mathcal{A}$} is a topos and
even when it is of the special type A\^{}, as I do not know of any \hand{$\mathcal{M}$} such
that
\[ (*) \qquad \hand{\mathcal{A}^{\mathcal{M}}\simeq\mathcal{A}}_{\mathrm{ab}} \qquad . \]
However, in case \hand{$\mathcal{A}$} = A\^{}, introducing the dual category B=A$^{\mathrm{o}}$ and \hand{$\mathcal{B}$} = B\^{},
we get (see D) below)
\[ (**) \qquad \hand{\mathcal{A}^{\mathcal{M}} \unsure{\cong|\approxeq} \mathcal{B}_{\mathcal{M}} \unsure{\cong|\approxeq} \mathcal{B}}_{\mathrm{ab}} \qquad \text{if} \quad \hand{\mathcal{M}} = \text{(Ab)} \quad , \]
hence we do get a canonical isomorphism \xout{provided} (*) provided A=B say
and hence \hand{$\mathcal{A}=\mathcal{B}$} . Let's look at any rate at the simplest case, namely
when A is \xout{the}\inserted{a} final object in (Cat), hence \hand{$\mathcal{A}$} can be identified with
(Sets), and \hand{$\mathcal{A}$}$_{\mathrm{ab}}$ with (Ab), the identification between the categories
\xout{(Ab) and Hom(} (Ab) and \uline{Hom}\hand{$_!$}(\xout{A}(Ab),(Ab)) being obtained \inserted{(we hope !)} by associating
to every object L in (Ab), the functor
\[ F \hand{\longmapsto} L\otimes F \; : \; \text{(Ab)} \hand{\longrightarrow} \text{(Ab)} \qquad . \]
This being so, the external operation (22) of (Ab) on \uline{Hom}\hand{$_!$} \hand{$\cong$} (Ab)
can be interpreted (using this identification) as the interior tensor
product operation in (Ab). On the other hand, the operation $*$ of (23)
is equally interpreted as nothing but the tensor product in (Ab), which
justifies the notation suggesting a tensor product. It would be nice
checking corresponding compatibilities for \xout{more}\inserted{\unsure{a}} general object\deleted{s} A in (Cat)
satisfying A=A$^{\mathrm{o}}$ , namely a direct sum of one-object cate\over{q}{g}ories A$_{\mathrm{i}}$ defi\xout{n}%
ned each in terms of a \uline{commutative} monoid M$_{\mathrm{i}}$ - I did'nt work it out
myself, sorry !

It should be noted that the relationship between the two \inserted{exterior} Hom's in
(29) and (31) is essentially the same as between the two exterior tensor-
type operations (26) and (23), Hom(F,K) designating an object in \hand{$\mathcal{N}$} and
\margin{$*$} \uline{Hom}(F,K) an \hand{$\mathcal{N}$}-valued sheaf on \hand{$\mathcal{A}$}, just as F \hand{$*$} L designates an object
\margin{$\circledast$} \xout{of} in \hand{$\mathcal{M}$} and F \hand{$\circledast$} L an \hand{$\mathcal{M}$}-valued cosheaf on \hand{$\mathcal{A}$}; the graphical device of
\uline{underlining} the symbol Hom (used for sheaves) corresponds to the device
\margin{$*$} of \uline{circling} the symbol $*$ (used for cosheaves). With this luxury of
explanations, I hope the notations introduced here are getting through...

\margin{75}
